% ============================================================
% CHAPTER 2: OBJECTIVES AND SCOPE
% ============================================================
\chapter{Objectives and Scope}

\section{Project Objectives}

The primary objective of the SafeGuard project is to design, develop, and deploy a comprehensive web-based personal safety management system. The specific objectives of the project are as follows:

\begin{enumerate}[label=\textbf{O\arabic*.}]
  \item \textbf{Design a Secure Authentication System:} To implement a robust user authentication mechanism using JWT (JSON Web Token) technology with bcrypt password hashing, ensuring secure access to the platform and protection of user credentials.

  \item \textbf{Develop an Emergency Contact Management Module:} To create a functional module that allows users to maintain a structured network of emergency contacts with associated metadata including relationship type, priority level, and multiple communication channels.

  \item \textbf{Implement a Real-Time SOS Alerting Mechanism:} To build an SOS emergency alert system that captures real-time geolocation data through the browser's Geolocation API and facilitates immediate notification of emergency contacts with customizable alert messages.

  \item \textbf{Create a Structured Incident Reporting System:} To develop an incident reporting module supporting multiple incident types, severity classification, geolocation tagging, evidence attachment, and lifecycle status tracking.

  \item \textbf{Build a Unified History and Analytics Dashboard:} To provide users with a comprehensive dashboard displaying aggregated statistics, recent activities, and historical records with filtering, search, and export capabilities.

  \item \textbf{Ensure Cross-Platform Accessibility:} To deliver a responsive web application accessible across devices and browsers, eliminating platform-specific dependencies through progressive enhancement.

  \item \textbf{Apply Modern Software Engineering Practices:} To demonstrate the application of industry-standard practices including RESTful API design, three-tier architecture, ORM-based data access, component-based UI development, and security best practices.
\end{enumerate}

\section{Scope}

\subsection{In Scope}

The scope of the SafeGuard project encompasses the following functional areas:

\begin{itemize}
  \item User registration, login, and profile management with JWT authentication
  \item Password change and profile picture upload functionality
  \item Emergency contact CRUD (Create, Read, Update, Delete) operations
  \item Contact search and priority-based filtering
  \item SOS alert creation with geolocation capture and multiple alert types
  \item SOS alert history with pagination, search, and status filtering
  \item Incident reporting with type classification, severity levels, and location data
  \item Incident status tracking through pending, investigating, resolved, and closed stages
  \item Incident image upload and attachment
  \item Combined alert and incident history with tabbed navigation
  \item CSV export of SOS alerts and incident records
  \item Statistics dashboard with aggregated metrics and recent activity feeds
  \item Activity logging for significant user actions
  \item Dark and light theme support with persistence
  \item Responsive design for desktop and mobile devices
  \item Role-based access control foundation (user and admin roles)
\end{itemize}

\subsection{Out of Scope}

The following features and functionalities are explicitly outside the scope of this project:

\begin{itemize}
  \item Real-time push notifications to emergency contacts via SMS or email
  \item Integration with third-party emergency services (police, ambulance, fire department)
  \item Video or voice calling capabilities
  \item Map-based visualization of alerts and incidents (GIS integration)
  \item Mobile native applications (iOS or Android)
  \item Multi-language (internationalization) support
  \item Offline data synchronization and caching
  \item Advanced analytics and machine learning-based threat detection
  \item Multi-factor authentication (MFA)
  \item Social media integration and sharing
\end{itemize}

\section{Applications}

The SafeGuard system has potential applications across several domains:

\begin{enumerate}[label=\arabic*.]
  \item \textbf{Personal Safety:} Individuals can use the platform to manage their personal safety, quickly alert trusted contacts during emergencies, and maintain a record of safety incidents.

  \item \textbf{Campus Security:} Educational institutions can deploy the system for students and staff, enabling centralized incident reporting and emergency communication.

  \item \textbf{Workplace Safety:} Organizations can utilize the platform for employee safety management, particularly for employees working in high-risk environments or field operations.

  \item \textbf{Travel Safety:} Travelers can use the system to maintain emergency contacts across locations and report incidents while travelling.

  \item \textbf{Community Safety:} Community organizations can leverage the platform for neighbourhood watch programs and community-based safety initiatives.
\end{enumerate}

\section{Benefits}

The SafeGuard project delivers benefits at multiple levels:

\subsection{User Benefits}
\begin{itemize}
  \item Single-platform solution for all personal safety needs
  \item Rapid emergency alerting with automatic location sharing
  \item Organized emergency contact management with priority classification
  \item Structured incident documentation for personal records and legal proceedings
  \item Accessibility from any device with a web browser
  \item Intuitive and aesthetically pleasing user interface
  \item Privacy through complete data isolation between users
\end{itemize}

\subsection{Technical Benefits}
\begin{itemize}
  \item Modern technology stack ensuring long-term maintainability
  \item RESTful API architecture enabling future integration with mobile applications
  \item Component-based frontend enabling rapid feature development
  \item Database normalization ensuring data integrity and efficient queries
  \item Security best practices including password hashing and token-based authentication
  \item Scalable architecture supporting horizontal and vertical scaling
\end{itemize}

\subsection{Academic Benefits}
\begin{itemize}
  \item Demonstrates application of full-stack web development concepts
  \item Applies database design and normalization principles
  \item Implements software engineering methodologies including requirements analysis, system design, and testing
  \item Provides practical experience with industry-standard frameworks and tools
\end{itemize}

\section{Limitations}

Despite its comprehensive feature set, the SafeGuard system has certain limitations:

\begin{enumerate}[label=\arabic*.]
  \item \textbf{Internet Dependency:} As a web application, SafeGuard requires an active internet connection to function. It does not support offline operation or data synchronization.

  \item \textbf{No Real-Time Notification Delivery:} The system records SOS alerts but does not actively push notifications to emergency contacts via SMS, email, or push notifications. Contacts must be manually notified.

  \item \textbf{No Native Mobile App:} The application is web-based and does not provide native mobile applications, which may limit accessibility in low-connectivity scenarios.

  \item \textbf{Geolocation Accuracy:} The accuracy of geolocation data depends on the device's GPS capabilities and may be limited indoors or in areas with poor satellite visibility.

  \item \textbf{Single-User Focus:} The current implementation focuses on individual personal safety and does not support group safety management or organizational deployment.

  \item \textbf{No Integration with Emergency Services:} The system does not integrate with police, ambulance, or fire department dispatch systems, limiting its utility in critical emergencies.

  \item \textbf{Limited Admin Functionality:} While the database schema supports admin roles, the current implementation provides limited administrative features for system management.

  \item \textbf{No Multi-Factor Authentication:} The authentication system relies solely on JWT tokens and does not implement multi-factor authentication for additional security.
\end{enumerate}
